\documentclass[11pt]{article}
\usepackage[margin=1in]{geometry}
\usepackage{amsmath,amssymb,amsthm}
\usepackage[T1]{fontenc}
\usepackage{lmodern}
\usepackage{microtype}
\usepackage{booktabs}
\usepackage[title]{appendix}
\usepackage[colorlinks=true,linkcolor=blue,citecolor=blue,urlcolor=blue]{hyperref}
\hypersetup{pdftitle={KPP fronts in the 1-1-1 ABC flow speed up linearly in every direction},pdfauthor={P. Graham}}
\newtheorem{theorem}{Theorem}
\newtheorem{proposition}[theorem]{Proposition}
\newtheorem{lemma}[theorem]{Lemma}
\theoremstyle{definition}
\newtheorem{definition}[theorem]{Definition}
\theoremstyle{remark}
\newtheorem{remark}[theorem]{Remark}
\newcommand{\R}{\mathbb{R}}
\newcommand{\Z}{\mathbb{Z}}
\newcommand{\C}{\mathbb{C}}
\newcommand{\T}{\mathbb{T}}
\newcommand{\Leb}{\operatorname{Leb}}
\title{KPP fronts in the 1-1-1 ABC flow speed up linearly in every direction}
\author{P. Graham\\ \small\url{https://orcid.org/0009-0007-8493-1505}}
\date{November 8, 2026}
\begin{document}
\maketitle
\begin{abstract}
Let $V(x,y,z)=(\sin z+\cos y,\ \sin x+\cos z,\ \sin y+\cos x)$ be the
Arnold--Beltrami--Childress (ABC) flow with $A=B=C=1$, and let $c^*_e(K)$ be
the minimal speed of pulsating fronts of $T_t+K\,V\cdot\nabla T=\Delta T+f(T)$
in the direction $e$, for a reaction $f$ of KPP type. We prove that
$\lim_{K\to\infty}c^*_e(K)/K$ exists and is positive for every unit vector
$e$. Zlatoš proved that this limit exists and equals a supremum over first
integrals of the flow, and Xin, Yu and Zlatoš, who constructed periodic orbits
of this flow, pointed out that a positive-measure tube of orbits moving
roughly linearly would make the limit positive, but that a rigorous proof of
such tubes was not known. We prove that such a tube exists. Two independent
interval-arithmetic computations certify an elliptic periodic orbit with
nonzero twist; Pöschel's KAM theorem with frequencies as parameters then gives
invariant curves of the return map with a Lipschitz label, which yield a
Lipschitz first integral supported on the tube. The deduction from the
certified facts is checked in Lean, with Zlatoš's and Pöschel's theorems as
cited inputs.

This work has not yet been reviewed by experts.
The proofs, the Lean formalization and the text were produced by AI models (from OpenAI and Anthropic), working under the author's direction.
\end{abstract}
\medskip
\noindent\textit{2020 Mathematics Subject Classification.} 35K57, 37J40, 76F25, 65G30.
\textit{Keywords.} ABC flow, KPP fronts, front speed-up, KAM theory, computer-assisted proof.

\section{Introduction}

\subsection{Main result}
Let $V\colon\R^3\to\R^3$ be the ABC flow with $A=B=C=1$,
\[
V(x,y,z)=(\sin z+\cos y,\ \sin x+\cos z,\ \sin y+\cos x).
\]
It is $2\pi$-periodic in each variable, divergence-free and of mean zero. Let
$f\colon[0,1]\to\R$ satisfy, for some $\varepsilon_0>0$,
\begin{equation}\label{eq:kpp}
\begin{gathered}
f\in C^{1,\varepsilon_0}([0,1]),\qquad f(0)=f(1)=0,\\
f \text{ non-increasing on } (1-\varepsilon_0,1),\qquad 0<f(s)\le f'(0)s \text{ for } s\in(0,1).
\end{gathered}
\end{equation}
For $K>0$ consider
\begin{equation}\label{eq:rd}
T_t+K\,V\cdot\nabla T=\Delta T+f(T)\qquad\text{on }\R\times\R^3.
\end{equation}
A \emph{pulsating front} in the direction of a unit vector $e$ with speed $c$
is a classical solution with values in $[0,1]$ of the form
$T(t,\xi)=U(\xi\cdot e-ct,\xi)$, where $U(s,\cdot)$ is $2\pi$-periodic for
every $s$, and $U(s,\xi)\to1$ as $s\to-\infty$ and $U(s,\xi)\to0$ as
$s\to+\infty$, uniformly in $\xi$. The minimal front speed $c^*_e(K)$ is the
infimum of the speeds of pulsating fronts; in this KPP setting fronts exist
exactly for $c\ge c^*_e(K)$ \cite{BH}.

\begin{theorem}\label{thm:main}
For every unit vector $e\in\R^3$ and every $f$ satisfying \eqref{eq:kpp}, the
limit $\lim_{K\to\infty}c^*_e(K)/K$ exists and is positive.
\end{theorem}

The proof uses a computer-assisted theorem about the streamlines of $V$
(Theorem~\ref{thm:orbit}), two published theorems, of Zlatoš
(Theorem~\ref{thm:zlatos}) and Pöschel (Theorem~\ref{thm:poschel}), quoted
exactly in Appendix~\ref{app:cited}, and a deduction that is checked in Lean
(Section~\ref{sec:lean}).

\subsection{The question and prior work}
Zlatoš \cite[Theorem~1.1]{Z} proved, for $1$-periodic flows $u$ and $f$ as in
\eqref{eq:kpp}, that $\lim_{A\to\infty}c^*_e(A)/A$ exists and equals a
supremum over the real first integrals $w\in H^1$ of $u$ (Theorem~\ref{thm:zlatos});
by \cite[Remark~3]{Z} it is positive ``precisely when there exists $w\in I$
such that $\int_{\T^n}(u\cdot e)w\,dx\ne0$''. El Smaily and Kirsch
\cite[Theorem~B and Remark~2]{ESK} extended the formula to periodic domains in
three or more dimensions and stated the same criterion. For the 1-1-1 ABC
flow, Xin, Yu and Zlatoš \cite{XYZ} proved that $V$ has periodic orbits with
rotation vectors parallel to the coordinate axes, which gives linear
speed-up for the G-equation model. In their introduction they wrote:
``numerical simulations suggest that there are `vortex tubes' composed of
orbits which travel roughly linearly in the $\pm x$, $\pm y$, and $\pm z$
directions [6], suggesting that $c^*(p,K)$ should also grow linearly in $K$
for each $p$. One could expect to find such vortex tubes, if they indeed
exist, in the vicinity of the `periodic' orbits constructed in this paper, but
a rigorous proof of their existence is currently not known.'' Here [6] is
Dombre, Frisch, Greene, Hénon, Mehr and Soward \cite{DFGHMS}. McMillen, Xin,
Yu and Zlatoš \cite[Section~1]{MXYZ}, working near the integrable regime
$0<A\ll1$, $B=C=1$, wrote: ``We are not able to prove the existence of such
tubes \dots\ even in the KAM region since KAM type theorems do not provide any
regularity of the set of quasi-periodic orbits.'' Shen, Xin and Zhou
\cite{SXZ} computed KPP front speeds numerically and found that ``in ABC
flows however, the enhancement is $O(\delta)$ which can be attributed to the
presence of principal vortex tubes''. Huang, Zhao and Dai \cite{HZD} study
invariant tori and chaotic streamlines in a near-integrable ABC regime,
using volume-preserving KAM and Melnikov methods.

\subsection{What is new}
We prove that a vortex tube exists, and we obtain the regularity that
\cite{MXYZ} asks for from a Lipschitz family of invariant curves rather than
from the set of quasi-periodic orbits itself. The proof has three new
ingredients. One is a computer-assisted proof that $V$ has an elliptic periodic orbit with nonzero
twist (Theorem~\ref{thm:orbit}). Breden and Lessard \cite{BL} gave
computer-assisted proofs for ABC flows: ``for ABC flows, we prove existence of
ballistic spiral orbits''. Kapela and Simó \cite{KS} proved rigorous KAM
results around periodic orbits, and Wilczak and Barrio \cite{WB} proved
stable elliptic periodic orbits with surrounding invariant tori, both by
computer. Another is the
passage from Pöschel's Theorem~A \cite{P}, a KAM theorem whose estimates are
explicit in all its parameters, to a positive-measure family of invariant
curves of the return map whose label is Lipschitz on their union
(Proposition~\ref{prop:kam}). This was the hardest step: theorems for a fixed
nondegenerate twist map do not survive the rescaling near an elliptic point,
and Pöschel's theorem applies only after a Floquet reduction, a time-periodic
Birkhoff normal form, and a choice of parameters that keeps the
perturbation below Pöschel's threshold. The last is a graph-form construction of
the first integral (Lemma~\ref{lem:w}). Recently Szumiński and Llibre
\cite{SzL} studied the ABC flow ``in perturbative regimes near its three
elementary integrable coordinate axes'' and proved the existence of elliptic
and hyperbolic periodic solutions there; their work concerns other parameter
values and other questions.

\section{Proof overview}
By Zlatoš's theorem, after rescaling to period $1$, the limit exists, and it
is positive as soon as $V$ has a Lipschitz, $2\pi$-periodic first integral $W$
with $\int (V\cdot e)W\ne0$. By the cyclic symmetry $(x,y,z)\mapsto(y,z,x)$ of
$V$ it suffices to find one $W\ge0$, supported where $V_z>0$, with
$\int V_xW=\int V_yW=0$ and $\int V_zW>0$. Such a $W$ comes from a tube of
orbits that move monotonically in $z$ and stay in a bounded region in $x$ and
$y$. We find the tube around an elliptic periodic orbit: its first-crossing
map $P$ from $\{z=0\}$ to $\{z=2\pi\}$ has an elliptic fixed point $p$
with nonzero twist (certified by computer), so KAM theory gives a
positive-measure family of $P$-invariant closed curves around $p$. The
measure of the set of labels of curves surrounding a point gives a
$P$-invariant Lipschitz function $w$ on the plane, and carrying $w$ along the
flow gives $W$. The one delicate point is regularity: $w$ is Lipschitz only
because the label of the curves is Lipschitz on their union, and that is what
Pöschel's estimates provide.

\section{The certified periodic orbit}\label{sec:orbit}
For $a\in\R^2$ let $P(a)\in\R^2$ be the point where the streamline of $V$
from $(a,0)$ first reaches the plane $z=2\pi$ (no reduction modulo $2\pi$).
Let $\mathcal X$ be the closed square with half-width $2^{-20}$ and center
\[
(6533234620664,\ 1727108826179)\cdot2^{-40}\approx(5.941942,\ 1.570796).
\]

\begin{theorem}[computer-assisted]\label{thm:orbit}
There is $p\in\mathcal X$ such that:
\begin{enumerate}
\item the streamline $\gamma$ of $V$ with $\gamma(0)=(p,0)$ satisfies
$\gamma(T)=(p,2\pi)$ for some $T\in[3.2352,3.2354]$, and
$\sin y+\cos x\ge1.92$ along $\gamma$ on $[0,T]$; in particular $P(p)=p$;
\item $P$ is differentiable at $p$ and $\operatorname{tr}DP(p)\in[-1.9650,-1.9640]$;
\item the first Birkhoff coefficient of $P$ at $p$ is nonzero.
\end{enumerate}
\end{theorem}

Since $P$ preserves the area form of Lemma~\ref{lem:ham} below and fixes $p$,
$\det DP(p)=1$; as the trace lies in
$(-2,-1)$, the eigenvalues of $DP(p)$ are $e^{\pm2\pi i\nu}$ with
$\nu\in(1/3,1/2)$, so $\lambda^q\ne1$ for $q=1,\dots,4$. The first Birkhoff
coefficient is defined in Appendix~\ref{app:twist} from the $3$-jet of $P$
at $p$; it is the coefficient of $w^2\bar w$ after the quadratic terms have
been removed in a complex eigencoordinate $w$, and its vanishing does not
depend on the coordinates.

Two programs prove Theorem~\ref{thm:orbit} independently. The first uses the
CAPD library \cite{CAPD}: it computes the $C^3$ jet of $P$ over all of $\mathcal X$,
proves a unique fixed point in $\mathcal X$ by the interval Newton method, encloses the
trace in $[-1.96455,-1.96451]$ and the return time in
$[3.2352771,3.2352775]$, and computes the Birkhoff normal form with CAPD's
routine, whose twist coefficient has real part in $[-0.5787,-0.5685]$. The
second, written independently, uses ball arithmetic from Arb \cite{Arb} with
its own Taylor integrator. It proves by the Krawczyk method a unique fixed
point in a square of half-width $2^{-100}$ inside $\mathcal X$, encloses the trace in
$-1.96453154969901637\pm4\cdot10^{-18}$ and $\det DP(p)$ in $1\pm3\cdot10^{-20}$,
and computes the twist coefficient of $\Phi\circ P\circ\Phi^{-1}$ for the
area-preserving coordinates $\Phi$ of Lemma~\ref{lem:ham}, with real part
$0.29530841922827626555718$, with error below $3\cdot10^{-24}$. Both check
$\sin y+\cos x>0$ on every integration step of the whole orbit segment, with
lower bounds $1.9215$ and $1.9378$. As a test, each program reproduces the
exact values for an integrable flow and rejects three altered certificates
(Appendix~\ref{app:checks}).
Appendix~\ref{app:checks} gives details.

\section{Invariant curves with a Lipschitz label}\label{sec:kam}
\begin{lemma}\label{lem:ham}
Let $X=x\sin y+\sin x$, $Y=y$, and
$H(x,y,z)=y\sin z+\sin y+\cos x-x\cos z$. Where $\sin y+\cos x\ne0$, the
streamlines of $V$, parametrized by $z$, satisfy
$dX/dz=\partial_YH$, $dY/dz=-\partial_XH$, with $H$ written in the
variables $(X,Y)$.
\end{lemma}

\begin{proof}
With $X$ fixed, $\partial x/\partial Y=-x\cos y/(\sin y+\cos x)$ and
$\partial x/\partial X=1/(\sin y+\cos x)$. Hence
$\partial_YH=V_x+x\cos y\,V_y/V_z=dX/dz$ and
$-\partial_XH=V_y/V_z=dY/dz$.
\end{proof}

So near the certified orbit the streamlines form a $2\pi$-periodic
Hamiltonian system with one degree of freedom, and $P$ is its period map,
area-preserving in $(X,Y)$.

\begin{proposition}\label{prop:kam}
Let $p$ be as in Theorem~\ref{thm:orbit} and $\delta>0$. There are numbers
$\rho,L>0$, a homeomorphism $h$ from a neighborhood of the closed disk
$\bar B(0,\rho)$ onto an open subset of $\R^2$ with $h(0)=p$,
$h(\bar B(0,\rho))\subset B(p,\delta)$ and $h$ bi-Lipschitz with constant
$L$ on $\bar B(0,\rho)$, a compact set $\mathcal K\subset\R$ of positive
measure, and continuous $2\pi$-periodic functions $g_s$, $s\in\mathcal K$,
with $0<g_s<\rho$, such that the closed curves
$C_s=h(\{g_s(\psi)(\cos\psi,\sin\psi):\psi\in\R\})$ satisfy
$P(C_s)=C_s$ and
\[
|s-s'|\le L\,\bigl|g_s(\psi)(\cos\psi,\sin\psi)-g_{s'}(\psi')(\cos\psi',\sin\psi')\bigr|
\quad\text{for all } s,s'\in\mathcal K,\ \psi,\psi'\in\R.
\]
\end{proposition}

\begin{proof}[Outline of proof]
The steps, each proved in Lean, are given in Appendix~\ref{app:kam}. Translate
to the periodic solution and pass to a real Floquet frame, in which the linear
part is the rotation with frequency $\nu$. Since $\lambda^q\ne1$ for $q\le4$, two
time-periodic symplectic changes of variables bring the Hamiltonian to
$\nu\,p+b\,p^2+R$ with $p=|w|^2/2$ and $R=O(|w|^5)$, analytic on a polydisc
times a strip $|\operatorname{Im}t|<s$. The twist coefficient is a
conjugacy invariant up to a positive factor, so Theorem~\ref{thm:orbit}(3)
gives $b\ne0$. For small $\varepsilon>0$ write
$p=p_0(\omega)+I$ with $p_0(\omega)=(\omega_1/\omega_2-\nu)/(2b)$ and treat
the frequencies $\omega=(\omega_1,\omega_2)$ as parameters, following
\cite[\S2a]{P}: the time-$t$ Hamiltonian becomes
$e(\omega)+\omega_1I+\omega_2J+P_\omega$ in action--angle variables
$(I,\theta_1,J,\theta_2=t)$, with $\omega$ in
$\Omega=\{|\omega_2-1|<\tfrac12,\ \omega_1/\omega_2\in\nu+2b(\varepsilon/2,\varepsilon)\}$.
Apply Theorem~\ref{thm:poschel} with $n=\tau=2$, $r=\varepsilon^{5/4}$,
$\alpha=\varepsilon^{9/8}$ and $h=\alpha s^3$: the perturbation is
$O(\varepsilon^{5/2})$, below Pöschel's threshold
$\gamma\alpha rs^3\sim\varepsilon^{19/8}$ for small $\varepsilon$. The time
row of the invariance equation forces $\omega_2=1$ to be preserved, so for
each $\sigma$ with $(\sigma,1)\in\Omega_\alpha$ the torus meets $\{t=0\}$ in
a $P$-invariant curve with rotation number $\sigma$. The labels $\sigma$ form
a compact set of measure at least
$|b|\varepsilon(1-8\alpha)-16\alpha(2+|\nu|)>0$. Pöschel's Lipschitz estimates
make each curve a radial graph over the angle and bound its displacement in
$\sigma$ below by $|\sigma-\sigma'|/(4|b|)$, minus terms of relative size
$O(\varepsilon^{1/8})$; this gives the Lipschitz label.
\end{proof}

\section{The invariant function and the tube}\label{sec:tube}
\begin{lemma}\label{lem:w}
In the setting of Proposition~\ref{prop:kam}, with $\delta$ small, let
$D_s=h(\{r(\cos\psi,\sin\psi):0\le r<g_s(\psi)\})$ and
$w(a)=\Leb\{s\in\mathcal K:a\in D_s\}$. Then $w\ge0$ is Lipschitz on $\R^2$,
$w(p)>0$, $w$ vanishes outside $h(B(0,\rho))$, and $w(P(a))=w(a)$ for $a$ in
$h(\bar B(0,\rho))$.
\end{lemma}

\begin{proof}
$w(p)=\Leb\mathcal K>0$, since $p=h(0)$ lies in every $D_s$. Let $a,b\in\R^2$.
If $a\in D_s$ and $b\notin D_s$ for some $s$, the segment between their
chart preimages meets the radial graph of $g_s$. So, for $a,b$ in
$h(\bar B(0,\rho))$, the labels $s$ counted in exactly one of $w(a)$, $w(b)$
are labels of points of one segment of length at most $L|a-b|$, and by the
Lipschitz label they form a set of diameter at most $L^2|a-b|$. Hence
$|w(a)-w(b)|\le L^2|a-b|$ there; since $w$ vanishes outside
$h(B(0,\rho))$, the bound extends to all of $\R^2$ with a larger constant. For invariance: near $p$, $P$ is a
homeomorphism onto its image (its inverse is the first-crossing map in
negative time), $P(C_s)=C_s$ and $P(p)=p$. So $P(D_s)$ is connected, misses
$C_s$ and contains $p$; as $h(\bar B(0,\rho))$ is a closed disk on which $h$
is a homeomorphism and $C_s$ is the boundary of $D_s$, this forces
$P(D_s)\subseteq D_s$, and the same argument for $P^{-1}$ gives equality.
\end{proof}

Radial graphs in a chart make the interior $D_s$ explicit, so no Jordan curve
theorem is needed.

\begin{proposition}\label{prop:tube}
There is a function $W\colon\R^3\to[0,\infty)$, not identically zero, that is
Lipschitz and $2\pi$-periodic in each variable, satisfies $V\cdot\nabla W=0$
almost everywhere, satisfies $V_z\ge c_0>0$ wherever $W\ne0$, and whose
support has $x$- and $y$-projections contained, modulo $2\pi$, in closed
intervals of length less than $2\pi$.
\end{proposition}

\begin{proof}
By Theorem~\ref{thm:orbit}(1) and Gronwall's inequality, $V_z\ge1.9$ along
every streamline from $(a,0)$ with $a$ near $p$ for $|t|\le4$, so each such
streamline crosses the planes $z=0$ and $z=\pm2\pi$ exactly once in that time
window, transversally. Choose $\delta$ in Proposition~\ref{prop:kam} small
enough that $w$ from Lemma~\ref{lem:w} is supported in this neighborhood.
For $\xi$ on the streamline through $(a,0)+\ell$, with $a$ in the support of
$w$ and $\ell\in2\pi\Z^3$, set $W(\xi)=w(a)$, and set $W=0$ elsewhere. This is
well defined, because the crossings of the streamline with the planes
$z\in2\pi\Z$ are the points $P^k(a)$ and $w\circ P=w$ there; it is constant
on streamlines; and it is Lipschitz, because the crossing point depends in a
Lipschitz way on $\xi$ (the flow is smooth and the crossings are
transversal). A Lipschitz function constant along streamlines satisfies
$V\cdot\nabla W=0$ wherever it is differentiable, hence almost everywhere.
On the support of $W$ we have $V_z=\sin y+\cos x\ge1.9$, which forces
$\cos x\ge0.9$ and $\sin y\ge0.9$; so, modulo $2\pi$, the $x$-projection of
the support lies in $[-\pi+1,\pi-1]$ and the $y$-projection in
$[-\pi/2+1,3\pi/2-1]$.
\end{proof}

\begin{lemma}\label{lem:avg}
For $W$ as in Proposition~\ref{prop:tube} and $Q=[0,2\pi)^3$,
$\int_QV_xW=\int_QV_yW=0$ and $\int_QV_zW>0$.
\end{lemma}

\begin{proof}
$\int_QV_zW\ge c_0\int_QW>0$. For $V_x$, let $\chi$ be a $2\pi$-periodic
Lipschitz function of $x$ alone that equals $x$ on the $x$-projection of the
support of $W$ (a sawtooth). Then $V_xW=(V\cdot\nabla\chi)W$, and, integrating
by parts over the torus with $\operatorname{div}V=0$,
$\int_Q(V\cdot\nabla\chi)W=-\int_Q\chi\,V\cdot\nabla W=0$. The same argument
applies to $V_y$.
\end{proof}

\section{Proof of Theorem~\ref{thm:main}}\label{sec:proof}
Let $u(\xi)=V(2\pi\xi)$ and $\hat f=(2\pi)^2f$. If $T$ is a front of
\eqref{eq:rd} with speed $c$, then $\hat T(t,\xi)=T((2\pi)^2t,2\pi\xi)$ is a
$1$-periodic front of $\hat T_t+2\pi K\,u\cdot\nabla\hat T=\Delta\hat T+\hat f(\hat T)$
with speed $2\pi c$, and conversely. Hence
$c^*_e(K)/K=\hat c^*_e(2\pi K)/(2\pi K)$, where $\hat c^*$ refers to $u$ and
$\hat f$. The flow $u$ satisfies Zlatoš's hypothesis (1.3) and $\hat f$ Zlatoš's
(1.2), so by Theorem~\ref{thm:zlatos} the limit exists and equals
$\ell(e)=\sup_w\int(u\cdot e)w^2/\|w\|_2^2$ over Zlatoš's set.

For positivity, let $\sigma$ be the cyclic permutation
$(x,y,z)\mapsto(y,z,x)$, which satisfies $V\circ\sigma=\sigma\circ V$, and let
$W$ be as in Proposition~\ref{prop:tube}. The functions $W$, $W\circ\sigma$
and $W\circ\sigma^{-1}$ are again Lipschitz periodic first integrals, and by
Lemma~\ref{lem:avg} and a change of variables,
$\int_Q(V\cdot e)W=e_z\int_QV_zW$,
$\int_Q(V\cdot e)\,W\circ\sigma=e_x\int_QV_zW$ and
$\int_Q(V\cdot e)\,W\circ\sigma^{-1}=e_y\int_QV_zW$. Choose the one whose
coefficient is a nonzero component of $e$, and let $\hat W(\xi)$ be it at
$2\pi\xi$. Then $\hat W$ lies in Zlatoš's set $I$ of real $H^1$ first
integrals of $u$, and $\beta=\int_{[0,1)^3}(u\cdot e)\hat W\ne0$. For small
$\eta>0$ the function $w_\eta=1+\eta\operatorname{sgn}(\beta)\hat W$ is in $I$,
satisfies $\|\nabla w_\eta\|_2^2=\eta^2\|\nabla\hat W\|_2^2\le\hat f'(0)\|w_\eta\|_2^2$,
and, since $\int u\cdot e=0$,
$\int(u\cdot e)w_\eta^2=2\eta|\beta|+\eta^2\int(u\cdot e)\hat W^2>0$. Hence
$\ell(e)>0$.\qed

\section{Formal verification and reproducibility}\label{sec:lean}
The deduction of Theorem~\ref{thm:main} from Theorem~\ref{thm:orbit} and the
two cited theorems is proved in Lean~4 \cite{Lean} with Mathlib
\cite{Mathlib}. The formal theorem is
\begin{quote}\raggedright
\texttt{theorem ABC.abc\_kpp\_linear\_speedup (f :} $\R\to\R$\texttt{) (hf : ZlatosKPP f) (e : E 3) (he :} $\|e\|=1$\texttt{) :}
$\exists\,\ell:\R,\ 0<\ell\ \wedge$ \texttt{Tendsto (fun K => cstar (2 * Real.pi) abcField f e K / K) atTop (}$\mathcal N\,\ell$\texttt{)}
\end{quote}
Here \texttt{abcField} is $V$; \texttt{ZlatosKPP f} is \eqref{eq:kpp} (an
exponent $\varepsilon\in(0,1]$, $f$ of class $C^1$ on $[0,1]$ with
$\varepsilon$-Hölder derivative, $f'(0)$ one-sided); and
\texttt{cstar (2 * Real.pi) abcField f e K} is the infimum of the speeds of
pulsating fronts as defined in Section~1, for classical solutions with $T$,
$T_t$, $\nabla T$ and $D^2T$ jointly continuous. The proof depends on
Lean's three standard axioms and on three further axioms:
\texttt{ABC.Cert.fixed\_point}, which states Theorem~\ref{thm:orbit} with its
numbers as exact rationals (the uniqueness of the fixed point, which both
programs also prove, is not used); \texttt{ABC.Zlatos2010.theorem\_1\_1}; and
\texttt{ABC.Poschel2001.theorem\_A}. The last two state
Theorems~\ref{thm:zlatos} and~\ref{thm:poschel} with all their hypotheses, and
the Lean files list how each phrase of the quoted statements is read (for
instance, in Pöschel's theorem, the sup-norms, ``real analytic'' as holomorphic and real
on real points, and the inclusions $\varphi(\Omega_\alpha)\subset O_h$, which
hold after decreasing its constant $\gamma$). Every hypothesis of the two
cited theorems is proved in Lean for this flow; nothing else is assumed. The
build checks the axiom list. A separate program checks that the rationals in the
certificate axiom are those the two programs certify.

What lies outside Lean is the correctness of the two interval computations,
including CAPD's normal-form routine and the Arb library, and the exact
transcription of the cited theorems, quoted in Appendix~\ref{app:cited} from
their arXiv versions (Zlatoš v2 and Pöschel's version 1.3, which notes small
corrections relative to the printed text). The commands to rebuild the Lean
proof and rerun both programs, with running times, are in the README of the
accompanying files.

\begin{appendices}
\section{The cited theorems}\label{app:cited}
\begin{theorem}[Zlatoš {\cite[Theorem~1.1]{Z}}]\label{thm:zlatos}
With (1.1) ``$T_t+Au\cdot\nabla T=\Delta T+f(T)$ on $\R^n$'', (1.2)
``$f\in C^{1,\varepsilon}([0,1])$, $f(0)=f(1)=0$ and $f$ is non-increasing on
$(1-\varepsilon,1)$ for some $\varepsilon>0$, $0<f(s)\le sf'(0)$ for
$s\in(0,1)$'', (1.3) ``The $1$-periodic flow $u\colon\T^n\to\R^n$ satisfies
$u\in C^{1,\varepsilon}(\T^n)$, $\nabla\cdot u\equiv0$, $\int_{\T^n}u\,dx=0$'',
and (1.5) ``$I\equiv\{w\in H^1(\T^n)\mid \operatorname{Im}w=0$ and
$u\cdot\nabla w=0\}$'': ``If $u$ and $f$ satisfy (1.2) and (1.3) and
$|e|=1$, then
\[
\lim_{A\to\infty}\frac{c^*_e(A)}{A}=\sup_{w\in I,\ \|\nabla w\|_2^2\le f'(0)\|w\|_2^2}
\frac{\int_{\T^n}(u\cdot e)w^2\,dx}{\|w\|_2^2}.
\]
In particular, the limit exists.''
\end{theorem}
Pulsating fronts are defined there as ``a solution of (1.1) of the form
$T(t,x)=U(x\cdot e-ct,x)$, with $c$ the front speed, and $U$ $1$-periodic in
$x$ and such that $\lim_{s\to-\infty}U(s,x)=1$,
$\lim_{s\to+\infty}U(s,x)=0$, uniformly in $x$''. In the formalization,
$u\cdot\nabla w=0$ holds almost everywhere, as in the proof of Zlatoš's
Theorem~1.2.

\begin{theorem}[Pöschel {\cite[Theorem~A]{P}}]\label{thm:poschel}
``Let $H=N+P$. Suppose $P$ is real analytic on $D_{r,s}\times O_h$ with
$|P|_{r,s,h}\le\gamma\alpha rs^\nu$, $\alpha s^\nu\le h$, where $\nu=\tau+1$
and $\gamma$ is a small constant depending only on $n$ and $\tau$. Suppose
also that $r,s,h\le1$. Then there exists a Lipschitz continuous map
$\varphi\colon\Omega_\alpha\to\Omega$ close to the identity and a Lipschitz
continuous family of real analytic torus embeddings
$\Phi\colon\T^n\times\Omega_\alpha\to B\times\T^n$ close to $\Phi_0$, such
that for each $\omega\in\Omega_\alpha$ the embedded tori are Lagrangian and
$X_H|_{\varphi(\omega)}\circ\Phi=D\Phi\cdot X_N$. Moreover, $\Phi$ is real
analytic on $T_*=\{\theta:|\operatorname{Im}\theta|<s/2\}$ for each $\omega$,
and
\[
|W(\Phi-\Phi_0)|,\ \alpha s^\nu|W(\Phi-\Phi_0)|_{\mathfrak L}\le\frac{c}{\alpha rs^\nu}|P|_{r,s,h},\qquad
|\varphi-\mathrm{id}|,\ \alpha s^\nu|\varphi-\mathrm{id}|_{\mathfrak L}\le\frac{c}{r}|P|_{r,s,h},
\]
uniformly on $T_*\times\Omega_\alpha$ and $\Omega_\alpha$, respectively, where
$c$ is a constant depending only on $n$ and $\tau$, and
$W=\operatorname{diag}(r^{-1}\mathrm{Id},s^{-1}\mathrm{Id})$.''
\end{theorem}
Here, in Pöschel's notation, $N=e(\omega)+\langle\omega,I\rangle$;
$\Omega_\alpha$ is the set of $\omega\in\Omega$ with
$|\langle k,\omega\rangle|\ge\alpha/|k|^\tau$ for all $0\ne k\in\Z^n$ and at
distance at least $\alpha$ from the boundary of $\Omega$, with $\tau>n-1$
fixed; $D_{r,s}=\{|I|<r\}\times\{|\operatorname{Im}\theta|<s\}$ and
$O_h=\{\omega:|\omega-\Omega_\alpha|<h\}$, with sup-norms; and
$\Phi_0(\theta,\omega)=(0,\theta)$. The Lean statement takes $n\ge2$, Pöschel's
standing assumption.

\section{The twist coefficient}\label{app:twist}
Let $F$ be a $C^3$ map of the plane with $F(p)=p$ and $\det DF(p)=1$, and let $v\in\C^2$ be an
eigenvector of $DF(p)$ for an eigenvalue $\lambda\notin\R$ with
$\lambda^2\ne\lambda$ and $\lambda^2\ne\bar\lambda$; then $|\lambda|=1$. In the coordinate $w$
given by $x=p+vw+\bar v\bar w$, write
$F=\sum_{j+k\le3}a_{jk}w^j\bar w^k+o(|w|^3)$, so $a_{10}=\lambda$,
$a_{01}=0$. The first Birkhoff coefficient is nonzero when
\[
\beta=a_{21}+\frac{2a_{20}a_{11}}{|\lambda|^2-\lambda}+\frac{a_{11}a_{20}}{\lambda^2-\lambda}
+\frac{|a_{11}|^2}{|\lambda|^2-\bar\lambda}+\frac{2|a_{02}|^2}{\lambda^2-\bar\lambda}\ne0,
\]
the coefficient of $w^2\bar w$ after the change $w\mapsto w+q(w)$ with
$q_{jk}=a_{jk}/(\lambda^j\bar\lambda^k-\lambda)$ that removes the quadratic
terms. The twist coefficient of both programs is $-i\beta/\lambda$ up to a
positive factor and complex conjugation; CAPD's routine uses only
$\bar\lambda=1/\lambda$, which holds at $p$ since $\det DP(p)=1$. Lean proves
that, under this determinant-one hypothesis, $\beta\ne0$ is invariant under
$C^3$ conjugacy fixing $p$, so the second
program's computation in the coordinates $\Phi$ of Lemma~\ref{lem:ham} gives
the same conclusion for $P$.

\section{The steps of Proposition~\ref{prop:kam}}\label{app:kam}
The Lean proof follows these steps; constants not given here are explicit in
the Lean files.
\begin{description}
\item[R1.] In the coordinates of Lemma~\ref{lem:ham}, translated by the
periodic solution, the streamlines solve $\dot u=J\nabla_uH_1(u,z)$ with
$H_1$ analytic on a polydisc times a strip, real, $2\pi$-periodic in $z$ and
$O(|u|^2)$. The period map $P_1$ is conjugate to $P$ by an analytic chart.
Analyticity of the flow in complex time and of $P_1$ is proved from
uniform limits of Euler polygons.
\item[R2.] With $X(t)$ the fundamental matrix of the linear part, the matrix
$Q(t)=X(t)SR(-\nu t)$, for $S\in SL(2,\R)$ bringing $DP_1(0)$ to the rotation
$R(2\pi\nu)$, is real, $2\pi$-periodic and symplectic, and $u=Q(t)w$ turns
$H_1$ into $\tfrac\nu2|w|^2+O(|w|^3)$. Here $\nu$ is not a multiple of $1/q$
for $q\le4$, by the certified trace.
\item[R3.] Two finite compositions of time-periodic polynomial shears, solving the homological
equations with divisors $m+\nu l$ bounded below for $|l|\le4$, give
$K=\tfrac\nu2|w|^2+\tfrac b4|w|^4+R$ with $|R|\le C|w|^5$ on a polydisc times
a strip, and conjugate the period maps.
\item[R4.] The period map of $K$ has twist coefficient a nonzero multiple of
$b$, and the conjugacy invariance of Appendix~\ref{app:twist} gives $b\ne0$.
\item[R5.] With $p=|w|^2/2=p_0(\omega)+I$ and polar angle $\theta_1$, the time
$\theta_2=t$ and its conjugate $J$, the Hamiltonian $\omega_2(K+J)$ equals
$e(\omega)+\omega_1I+\omega_2J+P_\omega$ exactly, with
$P_\omega=\omega_2\bigl(bI^2+R\bigr)$.
\item[R6.] For $\varepsilon\le\varepsilon_0$, with $r=\varepsilon^{5/4}$,
$\alpha=\varepsilon^{9/8}$ and $h=\alpha s^3$, $P_\omega$ is real analytic on
$D_{r,s}\times O_h$ with $|P_\omega|\le C_1\varepsilon^{5/2}$, and every
hypothesis of Theorem~\ref{thm:poschel} holds.
\item[R7.] The time row of $X_H\circ\Phi=D\Phi\cdot X_N$, averaged over the
torus, shows $\varphi_2(\omega)=\omega_2$ and that the $\theta_2$-component
of $\Phi-\Phi_0$ is constant, so for $(\sigma,1)\in\Omega_\alpha$ the section
$t=0$ of the torus is a $P$-invariant curve. The set
$\{\sigma:(\sigma,1)\in\Omega_\alpha\}$ has measure at least
$|b|\varepsilon(1-8\alpha)-16\alpha(2+|\nu|)$.
\item[Assembly.] With $\varepsilon=\kappa^8$ and $\kappa$ small, the angle
part of $\Phi$ minus the identity is $\tfrac12$-Lipschitz, so each curve is
the radial graph of $g_\sigma=\sqrt{2(p_0+\Phi^I)}$ over the angle, in the
chart given by R1--R3; the bounds on $\varphi-\mathrm{id}$ and
$\Phi-\Phi_0$ and on their Lipschitz constants give the Lipschitz label.
\end{description}

\section{The two programs}\label{app:checks}
The certificate lists the center and half-width of $\mathcal X$ and the claimed ranges
of the return time and the trace. The first program (CAPD 6.1.0, at a pinned
commit, with machine intervals) integrates the variational equations to
order $3$ over $\mathcal X$ by the $C^r$-Lohner algorithm \cite{WZ}, finds the crossing with $z=2\pi$, and checks every claim;
the second (Arb through python-flint 0.9.0, 256-bit precision, order-36
Taylor steps of size $1/32$ with a priori enclosures) does the same
independently, refining the fixed point in high precision first. As a test,
both treat the integrable flow $x'=\cos y$, $y'=\sin x$, $z'=13/10$ at $(0,\pi/2)$,
whose return map is the time-$2\pi/1.3$ map of $\cos x+\sin y$ with exact
twist $-(20\pi/13)/4$; both reproduce it (the second to within
$8\cdot10^{-29}$). Both reject three altered certificates: the center moved
by $2^{-19}$, the upper trace bound set to $-1.9646$, and the upper return
time set to $3.2352$.
\end{appendices}

\begin{thebibliography}{99}
\bibitem{BH} Henri Berestycki and François Hamel, Front propagation in
periodic excitable media, Comm. Pure Appl. Math. 55 (2002), 949--1032.
\bibitem{BL} Maxime Breden and Jean-Philippe Lessard, Polynomial interpolation
and a priori bootstrap for computer-assisted proofs in nonlinear ODEs,
Discrete Contin. Dyn. Syst. Ser. B 23 (2018), 2825--2858; arXiv:1704.03128.
\bibitem{CAPD} Tomasz Kapela, Marian Mrozek, Daniel Wilczak and Piotr
Zgliczyński, CAPD::DynSys: a flexible C++ toolbox for rigorous numerical
analysis of dynamical systems, Commun. Nonlinear Sci. Numer. Simul. 101
(2021), 105578.
\bibitem{DFGHMS} Thierry Dombre, Uriel Frisch, John M. Greene, Michel Hénon,
A. Mehr and Andrew M. Soward, Chaotic streamlines in the ABC flows, J. Fluid
Mech. 167 (1986), 353--391.
\bibitem{ESK} Mohammad El Smaily and Stéphane Kirsch, Front speed enhancement
by incompressible flows in three or higher dimensions, Arch. Ration. Mech.
Anal. 213 (2014), 327--354; arXiv:1307.4106v2.
\bibitem{HZD} De-Bin Huang, Xiao-Hua Zhao and Hui-Hui Dai, Invariant tori and
chaotic streamlines in the ABC flow, Phys. Lett. A 237 (1998), 136--140.
\bibitem{Arb} Fredrik Johansson, Arb: efficient arbitrary-precision
midpoint-radius interval arithmetic, IEEE Trans. Comput. 66 (2017),
1281--1292.
\bibitem{KS} Tomasz Kapela and Carles Simó, Rigorous KAM results around
arbitrary periodic orbits for Hamiltonian systems, Nonlinearity 30 (2017),
965--986.
\bibitem{Mathlib} The mathlib Community, The Lean mathematical library, in:
Proceedings of the 9th ACM SIGPLAN International Conference on Certified
Programs and Proofs (CPP 2020), 367--381.
\bibitem{MXYZ} Tyler McMillen, Jack Xin, Yifeng Yu and Andrej Zlatoš,
Ballistic orbits and front speed enhancement for ABC flows, SIAM J. Appl.
Dyn. Syst. 15 (2016), 1753--1782; arXiv:1606.01461.
\bibitem{Lean} Leonardo de Moura and Sebastian Ullrich, The Lean~4 theorem
prover and programming language, in: Automated Deduction (CADE 28), Lecture
Notes in Computer Science 12699, Springer, 2021, 625--635.
\bibitem{P} Jürgen Pöschel, A lecture on the classical KAM theorem, in:
Smooth Ergodic Theory and Its Applications, Proc. Sympos. Pure Math. 69
(2001), 707--732; arXiv:0908.2234v1 (version 1.3, 2009).
\bibitem{SXZ} Lyonell Shen, Jack Xin and Aimin Zhou, Finite element
computation of KPP front speeds in 3D cellular and ABC flows, Math. Model.
Nat. Phenom. 8 (2013), no.~3, 182--197.
\bibitem{SzL} Wojciech Szumiński and Jaume Llibre, Dynamics, periodic orbits
and $C^1$ non-integrability of the ABC flow, arXiv:2609.17449 (2026).
\bibitem{WB} Daniel Wilczak and Roberto Barrio, Systematic computer-assisted
proof of branches of stable elliptic periodic orbits and surrounding invariant
tori, SIAM J. Appl. Dyn. Syst. 16 (2017), 1618--1649.
\bibitem{WZ} Daniel Wilczak and Piotr Zgliczyński, $C^r$-Lohner algorithm,
arXiv:0704.0720 (2007).
\bibitem{XYZ} Jack Xin, Yifeng Yu and Andrej Zlatoš, Periodic orbits of the
ABC flow with $A=B=C=1$, SIAM J. Math. Anal. 48 (2016), 4087--4093;
arXiv:1601.02724.
\bibitem{Z} Andrej Zlatoš, Sharp asymptotics for KPP pulsating front speed-up
and diffusion enhancement by flows, Arch. Ration. Mech. Anal. 195 (2010),
441--453; arXiv:0704.1163v2.
\end{thebibliography}
\end{document}
